% =============================================================
% Appendix A.7: The General Geometric Model.
%
% Source: ~/evolmoves/doc/tex/main.tex, subsection "The GGI and
% TKF92 models" (lines ~324-426 of that file at the time of this
% paper's release). Quoted near-verbatim with the following local
% adaptations:
%   * \fragext -> \ext (the supplement uses \ext for the TKF92
%     fragment-extension probability r).
%   * \prob{k}     -> P(k)
%     \prob{n}     -> P(n)
%   * Stripped \index{...} commands.
%   * Removed a forward-reference to ``Section tkftriad'' that does
%     not exist in this supplement (the corresponding triad
%     construction lives in the Maraschino / VarAnc subsections of
%     Appendix B and is not the relevant cross-reference here).
%
% The \section{} below becomes a \subsection under \demote in the
% outer supplement, i.e. it renders as A.7.
% =============================================================

\section{The General Geometric model}
\label{sec:ggi-tkf92}

Ideally, the underlying evolutionary model we'd use for indels would be the General Geometric Indel (GGI) model \cite{Holmes2020,DeMaio2020}.
The GGI model is, in some sense, the maximum-entropy indel model for a given
expected rate and length of indels.
Unfortunately, the probability distributions of aligned ancestor-descendant sequences
at finite times under this model are hard to solve.
It turns out that the pairwise alignment distributions under the GGI model are well approximated by the TKF92 model \cite{ThorneEtAl92}, so we use that model for inference.
In this section we discuss the relationship between the models
and, in particular, how we can map parameters between them in a principled way.

The GGI model has parameters $(\insrate_0,\delrate_0,x,y)$.
In this model, the instantaneous rate of insertion of $k$
residues at any available site is $\insrate_0 x^{k-1} (1-x)$
(with the residues themselves sampled from a stationary distribution),
while the instantaneous rate of deletion of a stretch of exactly $k$ contiguous residues
is $\delrate_0 y^{k-1} (1-y)$.
For reversibility (or, indeed, the existence of an equilibrium) we require that
$\insrate_0 y (1-x) = \delrate_0 x (1-y)$.
The mean insertion and deletion lengths are $1/(1-x)$ and $1/(1-y)$.
The rate at which any given residue is deleted,
counting deletions that start at or before that residue (in an infinitely long sequence),
is $\delrate_0 / (1-y)$.
The probability that a sequence has length $n$ at equilibrium is $(x/y)^n(1-x/y)$,
so the expected sequence length at equilibrium
is $\ell_{\mathrm{GGI}} = \frac{x/y}{1-x/y} = \frac{1}{y/x-1}$.

The TKF92 model has parameters $(\insrate,\delrate,\ext)$.
In this model,
$k$-residue fragments (where $k>0$ and $P(k) = \ext^{k-1}(1-\ext)$) are inserted and deleted at rates $\insrate,\delrate$.
The rate at which any given residue is deleted
is just the rate at which its fragment is deleted, $\delrate$.
The mean number of fragments at equilibrium is
$\frac{\insrate/\delrate}{1-\insrate/\delrate} = \frac{1}{\delrate/\insrate - 1}$
and the mean fragment length is $1/(1-\ext)$,
so the mean sequence length at equilibrium is
$\ell_{\mathrm{TKF}} = \frac{1}{(\delrate/\insrate - 1)(1 - \ext)}$.
The precise distribution of sequence lengths at equilibrium is
a zero-altered geometric distribution
\[
P(n) = \left\{
\begin{array}{ll}
1 - \frac{\insrate}{\delrate} & \mbox{if $n=0$} \\
\frac{\insrate}{\delrate} \left( \ext + (1-\ext)\frac{\insrate}{\delrate} \right)^{n-1} (1-\ext) \left( 1 - \frac{\insrate}{\delrate} \right) & \mbox{if $n>0$}
\end{array}
\right.
\]

The mean length of an insertion in TKF92 is $1/(1-\ext)$, i.e.\ the expected fragment length.
For deletions, things are slightly more complicated. Superficially,
it looks like deletions are also fragments so they should have the same mean length.
However, if we do not have knowledge of the fragment boundaries,
but are conditioning just on the current sequence length (as in GGI),
we have to weight this by the posterior probability that the subsequence being deleted does in fact constitute a single fragment.
For a subsequence of length $k$, this posterior probability is proportional to
$\left(\frac{\ext}{\ext + (1-\ext)\frac{\insrate}{\delrate}}\right)^{k-1}$
(indels at the end of the sequence have a slightly different correction factor
since they are more likely to constitute a fragment, but the $k$-dependence is the same).
The rate of starting an insertion or deletion {\em at a particular site}
is subject to a similar correction factor.

Given the equilibrium constraint on GGI, each model has three free parameters.
In order to find a correspondence, it is useful to match moments (expectations).
In view of the above discussion, the most natural expectations to equate are
(in the form ``GGI expectation = TKF92 expectation'')
\begin{itemize}
    \item The mean deletion rate of a given residue: $\delrate_0/(1-y) = \delrate$
    \item The mean length at equilibrium: $\frac{1}{y/x-1} = \frac{1}{(\delrate/\insrate - 1)(1 - \ext)}$
    \item The mean length of an insertion event: $1/(1-x) = 1/(1-\ext)$ so $x=\ext$.
\end{itemize}

These equations, together with the reversibility constraint
$\insrate_0 y (1-x) = \delrate_0 x (1-y)$,
fully specify a one-to-one mapping between GGI and TKF92.
In what follows, we will simply use the TKF92 parameters
$(\insrate, \delrate, \ext)$ and likelihoods
(including the zero-adjusted equilibrium length distribution),
but with the understanding that we are approximating GGI.
In particular, when we jointly sample the pairwise alignments around any
given node in the phylogenetic tree, we will {\em not} require the fragment
boundaries to coincide (as would be the case in TKF92).
Technically speaking, any indel model parameter priors should be on the GGI parameters
rather than the TKF92 ones, but since we are using uninformative priors
this may be an acceptable compromise.

One way in which the TKF92 model, where the process state is augmented with latent fragment boundaries,
deviates from the sequence-only GGI is that when we use the TKF92 Pair HMM
(with the fragment boundaries marginalized out) as an approximation to $P(X_t|X_0)$,
we cannot expect it perfectly to obey the
Chapman-Kolmogorov equation $P(X_{t+u}=x''|X_0=x) = \sum_{x'} P(X_t=x'|X_0=x)P(X_{t+u}=x''|X_t=x')$.
This has consequences for any MCMC moves that split branches using a Pair HMM
(which models exactly the Chapman-Kolmogorov equation), since it means that those moves will be
miscalibrated in probability, but the approximation should be mild.
In fact, it seems possible that no five-state ($\smide$) Pair HMM exactly obeys the Chapman-Kolmogorov equation
for a GGI model, except in the special case where the model reduces to the TKF91 model.
The best approximations known are obtained by matching moments of the underlying counting process
and thereby deriving ODEs for the transition probabilities \cite{Holmes2020}.
However, the TKF92 model is a close runner-up (and actually appears to fit real data better).
